\section{Introduction}
\subsection{Presentation of the host laboratories}
\myindent This intership wac conducted as part of my Master 2 at the Magistère of Fundamental Physics of Paris-Saclay Université, which I 
completed at Aix-Marseille Université in the ``Physique Fondamentale et Applications'' track. It was carried out at the 
``Physique des Cellules et Cancer'' (PCC, UMR 168) laboratory of the Institut Curie as well as at the Gulliver laboratory (UMR 7083) 
of the ``École Supérieure de Physique et de Chimie Industrielles'' (ESPCI) in Paris, and was supervised by Michele Castellana 
(Institut Curie) and David Lacoste (ESPCI). Although the PCC unit is a leading biological physics laboratory, part of the 
Institut Curie -- a pioneer institution in cancer research and treatment -- it host a reknowned theory team (of which Michele Castellana 
is a member) that is sometimes interested in physics-driven topics such as the one invastigated in this report. The Gulliver
 laboratory hosts both experimentalists and theorician (among whom David Lacoste), working at the frontieer between physical, 
 chemical and biological topics. 

\subsection{Subject}
\myindent  In October 2024, John J.Hopfield was awarded the Nobel Prize in Physics alongside Geoffrey Hinton \cite{nobel2024physics} for their pioneer work on artifical neural 
networks. This work laid the foundations for further development and understanding of neural networks, whether they are organic or artificial, and eventually lead to 
the arising of Artificial Intelligience (AI).\\
Hopfield worked on a artificial neural network, originally invented by Shun'ichi Amari in 1972 \cite{Amari1972LearningPA}. He greatly contributed to the analysis, 
understanding and popularization of the network in his landmark paper \textit{Neural networks and physical systems with emergent collective computational abilities} 
published in 1982, and the network was eventually named after him.
\cite{Hopfield1982}
\subsection{Objectives}
